% !TEX root = ../projeto-pesquisa.tex
% =====================================================================
%  RESUMO
% =====================================================================
\secao{Resumo}
\justifying
O avanço das tecnologias e dos novos modelos generativos de inteligência
artificial tem apresentado um cenário desafiador no que tange à identificação de
imagens falsas, cuja sofisticação amplia o potencial de disseminação de
desinformação e golpes digitais no ambiente \textit{online}. Nesse contexto, o
presente projeto propõe o desenvolvimento de solução tecnológica voltada à
análise de imagens e capaz de apoiar a verificação da autenticidade e/ou
manipulação de arquivos digitais, tendo por objetivos: a) a implementação de
algoritmos de extração e análise de metadados e de padrões de compressão para
identificar assinaturas digitais características de processos de edição por IA;
b) a avaliação do potencial impacto do \textit{software} como mecanismo de
validação de conteúdo multimídia em larga escala na mitigação de desinformação;
c) a realização de testes comparativos de acurácia entre o \textit{software}
desenvolvido e o discernimento humano, estabelecendo métricas de desempenho para
a identificação de imagens sintéticas e d) a investigação das limitações da
técnica de \textit{Error Level Analysis} (ELA) frente ao refinamento de modelos
generativos, verificando a ocorrência de padronização de ruídos e compressão em
imagens de alta fidelidade. Para tanto, será realizada ação de extensão com cerca
de 40 estudantes de ensino médio de instituição de ensino em Santo André --- cuja
parceria ainda será formalizada, e que consistirá em atividade experimental com
os alunos, que divididos em dois grupos avaliarão cerca de 20 imagens
padronizadas (10 reais e 10 artificiais), sem e com o apoio da ferramenta
tecnológica desenvolvida para o projeto. Espera-se obter uma análise estatística
dos resultados com a identificação de padrões de erro, visando avaliar a eficácia
das técnicas computacionais frente à percepção humana, além de promover a
reflexão crítica sobre desinformação digital e o uso consciente de tecnologias.

\vspace{5mm}
\noindent\textbf{Palavras-chave:} análise de imagens, ciência de dados,
\textit{deepfakes}, desinformação, inteligência artificial.

% =====================================================================
%  ABSTRACT
% =====================================================================
\newpage
\secao{Abstract}
\justifying
\begin{otherlanguage}{english}
The advancement of technologies and new generative artificial intelligence models
has presented a challenging scenario regarding the identification of fake images,
whose sophistication increases the potential for the dissemination of
misinformation and digital scams in the online environment. In this context, this
project proposes the development of a technological solution focused on image
analysis and capable of supporting the verification of the authenticity and/or
manipulation of digital files, with the following objectives: a) the
implementation of metadata extraction and analysis algorithms and compression
patterns to identify digital signatures characteristic of AI editing processes;
b) the evaluation of the potential impact of the software as a mechanism for
large-scale multimedia content validation in mitigating misinformation; c) the
performance of comparative accuracy tests between the developed software and
human discernment, establishing performance metrics for the identification of
synthetic images; and d) the investigation of the limitations of the Error Level
Analysis (ELA) technique in the face of the refinement of generative models,
verifying the occurrence of noise standardization and compression in
high-fidelity images. To this end, an outreach activity will be carried out with
approximately 40 high school students from an educational institution in Santo
André --- whose partnership is yet to be formalized --- and will consist of an
experimental activity with the students, who, divided into two groups, will
evaluate approximately 20 standardized images (10 real and 10 artificial), both
with and without the support of the technological tool developed for the project.
The aim is to obtain a statistical analysis of the results, identifying error
patterns, in order to evaluate the effectiveness of computational techniques
compared to human perception, as well as promoting critical reflection on digital
misinformation and the conscious use of technology.

\vspace{5mm}
\noindent\textbf{Keywords:} artificial intelligence, data science, deepfakes,
disinformation, image analysis.
\end{otherlanguage}

